\chapter{Calculus}
\label{\detokenize{tools-for-mathematics/03-calculus/introduction/main:calculus}}\label{\detokenize{tools-for-mathematics/03-calculus/introduction/main::doc}}

The A-level syllabus describes Calculus 
describes
as:

\begin{quote}
``Calculus: this is a fundamental element which describes change in dynamic
situations and underlines the links between functions and graphs.''
\end{quote}


In practice this often means:
\begin{itemize}
\item 

taking \index{limits} of functions;

\item 

\index{differentiating} functions;

\item 

\index{integrating} functions.

\end{itemize}


Here you will see how to instruct a computer to carry out these techniques.



\begin{note}
In this chapter you will cover:
\begin{itemize}
\item 

Taking \index{limits}.

\item 

\index{differentiating} functions.

\item 

Computing definite and indefinite integrals.

\end{itemize}
\end{note}





\section{Tutorial}
\label{\detokenize{tools-for-mathematics/03-calculus/tutorial/main:tutorial}}\label{\detokenize{tools-for-mathematics/03-calculus/tutorial/main::doc}}

You will solve the following problem using a computer to assist with the
technical aspects:

Consider the function \(f(x)= \frac{24 x \left(a - 4 x\right) + 2 \left(a - 8 x\right) \left(b - 4 x\right)}{\left(b - 4 x\right)^{4}}\)
\begin{enumerate}

\item 

Given that \(\frac{df}{dx}|_{x=0}=0\), \(\frac{d^2f}{dx^2}|_{x=0}=-1\) and that \(b>0\) find the values of \(a\) and \(b\).

\item 

For the specific values of \(a\) and \(b\) find:
\begin{enumerate}

\item 

\(\lim_{x\to 0}f(x)\);

\item 

\(\lim_{x\to \infty}f(x)\);

\item 

\(\int f(x) dx\);

\item 

\(\int_{5}^{20} f(x) dx\).

\end{enumerate}

\end{enumerate}



Sympy is once again the library you will use for this.
You will start by creating a variable \texttt{expression} that has the value of the expression of \(f(x)\):




\begin{pyin}
import sympy as sym

x = sym.Symbol("x")
a = sym.Symbol("a")
b = sym.Symbol("b")
expression = (24 * x * (a - 4 * x) + 2 * (a - 8 * x) * (b - 4 * x)) / ((b - 4 * x) ** 4)
expression
\end{pyin}




\begin{equation*}
\begin{split}\displaystyle \frac{24 x \left(a - 4 x\right) + \left(2 a - 16 x\right) \left(b - 4 x\right)}{\left(b - 4 x\right)^{4}}\end{split}
\end{equation*}

You will use \texttt{sympy.diff} to calculate the derivative. This tool takes two inputs:
\begin{itemize}
\item 

the first is the expression you are \index{differentiating}. Essentially this is the numerator of \(\frac{df}{dx}\).

\item 

the second is the variable you are \index{differentiating} with respect to. This is the denominator of \(\frac{df}{dx}\).

\end{itemize}



You have imported \texttt{import sympy as sym} so you are going to write \texttt{sym.diff}:





\begin{pyin}
derivative = sym.diff(expression, x)
derivative
\end{pyin}




\begin{equation*}
\begin{split}\displaystyle \frac{16 a - 16 b - 64 x}{\left(b - 4 x\right)^{4}} + \frac{16 \cdot \left(24 x \left(a - 4 x\right) + \left(2 a - 16 x\right) \left(b - 4 x\right)\right)}{\left(b - 4 x\right)^{5}}\end{split}
\end{equation*}




Factorise that to make it slightly clearer:




\begin{pyin}
sym.factor(derivative)
\end{pyin}




\begin{equation*}
\begin{split}\displaystyle \frac{16 \left(- 3 a b - 12 a x + b^{2} + 16 b x + 16 x^{2}\right)}{\left(- b + 4 x\right)^{5}}\end{split}
\end{equation*}

You will now create the first equation, which is obtained by substituting \(x=0\)
in to the value of the derivative and equating that to \(0\):




\begin{pyin}
first_equation = sym.Eq(derivative.subs({x: 0}), 0)
first_equation
\end{pyin}




\begin{equation*}
\begin{split}\displaystyle \frac{32 a}{b^{4}} + \frac{16 a - 16 b}{b^{4}} = 0\end{split}
\end{equation*}




Factor that equation:




\begin{pyin}
sym.factor(first_equation)
\end{pyin}




\begin{equation*}
\begin{split}\displaystyle \frac{16 \cdot \left(3 a - b\right)}{b^{4}} = 0\end{split}
\end{equation*}




Now you are going to create the second equation, substituting \(x=0\) in to the
value of the second derivative. Calculate the second derivative by passing a
third (optional) input to \texttt{sym.diff}:




\begin{pyin}
second_derivative = sym.diff(expression, x, 2)
second_derivative
\end{pyin}




\begin{equation*}
\begin{split}\displaystyle \frac{64 \left(-1 - \frac{8 \left(- a + b + 4 x\right)}{b - 4 x} + \frac{10 \cdot \left(12 x \left(a - 4 x\right) + \left(a - 8 x\right) \left(b - 4 x\right)\right)}{\left(b - 4 x\right)^{2}}\right)}{\left(b - 4 x\right)^{4}}\end{split}
\end{equation*}

Equate this expression to \(-1\):

\begin{pyin}
second_equation = sym.Eq(second_derivative.subs({x: 0}), -1)
second_equation
\end{pyin}




\begin{equation*}
\begin{split}\displaystyle \frac{64 \cdot \left(\frac{10 a}{b} - 1 - \frac{8 \left(- a + b\right)}{b}\right)}{b^{4}} = -1\end{split}
\end{equation*}




Now solve the first equation to obtain a value for \(a\):




\begin{pyin}
sym.solveset(first_equation, a)
\end{pyin}




\begin{equation*}
\begin{split}\displaystyle \left\{\frac{b}{3}\right\}\end{split}
\end{equation*}




Now to substitute that value for \(a\) and solve the second equation for \(b\):




\begin{pyin}
second_equation = second_equation.subs({a: b / 3})
second_equation
\end{pyin}




\begin{equation*}
\begin{split}\displaystyle - \frac{192}{b^{4}} = -1\end{split}
\end{equation*}






\begin{pyin}
sym.solveset(second_equation, b)
\end{pyin}




\begin{equation*}
\begin{split}\displaystyle \left\{- 2 \sqrt{2} \cdot \sqrt[4]{3}, 2 \sqrt{2} \cdot \sqrt[4]{3}, - 2 \sqrt{2} \cdot \sqrt[4]{3} i, 2 \sqrt{2} \cdot \sqrt[4]{3} i\right\}\end{split}
\end{equation*}




Recalling the question you know that \(b>0\) thus: \(b = 2\sqrt{2}\sqrt[4]{3}\) and
\(a=\frac{2\sqrt{2}\sqrt[4]{3}}{3}\).
You will substitute these values back and finish the question:


\begin{pyin}
expression = expression.subs(
    {
        a: 2 * sym.sqrt(2) * sym.root(3, 4) / 3, 
        b: 2 * sym.sqrt(2) * sym.root(3, 4),
    }
)
expression
\end{pyin}




\begin{equation*}
\begin{split}\displaystyle \frac{24 x \left(- 4 x + \frac{2 \sqrt{2} \cdot \sqrt[4]{3}}{3}\right) + \left(- 16 x + \frac{4 \sqrt{2} \cdot \sqrt[4]{3}}{3}\right) \left(- 4 x + 2 \sqrt{2} \cdot \sqrt[4]{3}\right)}{\left(- 4 x + 2 \sqrt{2} \cdot \sqrt[4]{3}\right)^{4}}\end{split}
\end{equation*}

You are using the \texttt{sym.root} command for the generic \(n\)th root.
You can confirm this:




\begin{pyin}
sym.diff(expression, x).subs({x: 0})
\end{pyin}




\begin{equation*}
\begin{split}\displaystyle 0\end{split}
\end{equation*}






\begin{pyin}
sym.diff(expression, x, 2).subs({x: 0})
\end{pyin}




\begin{equation*}
\begin{split}\displaystyle -1\end{split}
\end{equation*}




Now you will calculate the \index{limits} using \texttt{sym.limit}, this takes 3 inputs:
\begin{itemize}
\item 

The expression you are taking the limit of.

\item 

The variable that is changing.

\item 

The value that the variable is tending towards.

\end{itemize}




\begin{pyin}
sym.limit(expression, x, 0)
\end{pyin}




\begin{equation*}
\begin{split}\displaystyle \frac{\sqrt{3}}{36}\end{split}
\end{equation*}






\begin{pyin}
sym.limit(expression, x, sym.oo)
\end{pyin}




\begin{equation*}
\begin{split}\displaystyle 0\end{split}
\end{equation*}




Now you are going to calculate the \textbf{indefinite} integral using
\texttt{sympy.integrate}. This tool takes 2 inputs as:
\begin{itemize}
\item 

the first is the expression you’re \index{integrating}. This is the \(f\) in \(\int_a^b f
dx\).

\item 

the second is the remaining information needed to calculate the integral: \(x\).

\end{itemize}




\begin{pyin}
sym.factor(sym.integrate(expression, x))
\end{pyin}




\begin{equation*}
\begin{split}\displaystyle \frac{x \left(6 x - \sqrt{2} \cdot \sqrt[4]{3}\right)}{12 \cdot \left(4 x^{3} - 6 \sqrt{2} \cdot \sqrt[4]{3} x^{2} + 6 \sqrt{3} x - \sqrt{2} \cdot 3^{\frac{3}{4}}\right)}\end{split}
\end{equation*}




If you want to calculate a \textbf{definite} integral then instead of passing the
single variable you pass a tuple which contains the variables as the
bounds of integration:




\begin{pyin}
sym.factor(sym.integrate(expression, (x, 5, 20)))
\end{pyin}




\begin{equation*}
\begin{split}\displaystyle - \frac{5 \left(- 5000 \sqrt{2} \cdot \sqrt[4]{3} - 1200 \sqrt{3} + 75 \sqrt{2} \cdot 3^{\frac{3}{4}} + 119997\right)}{2 \left(-32000 - 120 \sqrt{3} + \sqrt{2} \cdot 3^{\frac{3}{4}} + 2400 \sqrt{2} \cdot \sqrt[4]{3}\right) \left(-500 - 30 \sqrt{3} + \sqrt{2} \cdot 3^{\frac{3}{4}} + 150 \sqrt{2} \cdot \sqrt[4]{3}\right)}\end{split}
\end{equation*}



\begin{note}
This tutorial has:
\begin{itemize}

\item 
    Simplified a rational quotient;

\item 

Differentiated symbolic expressions;


\item 

Solved algebraic equations.

\end{itemize}
\end{note}




\section{How to}

\subsection{Calculate the derivative of an expression.}

You can calculate the derivative of an expression using \texttt{sympy.diff} which takes,
an expression, a variable and a degree.


\begin{api}
sympy.diff(expression, variable, degree=1)
\end{api}



The default value of \texttt{degree} is 1.
For example to compute \(\frac{d (4 x ^ 3 + 2 x + 1}{dx}\):




\begin{pyin}
import sympy as sym

x = sym.Symbol("x")
expression = 4 * x ** 3 + 2 * x + 1
sym.diff(expression, x)
\end{pyin}




\begin{equation*}
\begin{split}\displaystyle 12 x^{2} + 2\end{split}
\end{equation*}




To compute the second derivative: \(\frac{d ^ 2 (4 x ^ 3 + 2 x + 1}{dx ^ 2}\)




\begin{pyin}
sym.diff(expression, x, 2)
\end{pyin}




\begin{equation*}
\begin{split}\displaystyle 24 x\end{split}
\end{equation*}




\subsection{Calculate the indefinite integral of an expression.}

You can calculate the indefinite integral of an expression using
\texttt{sympy.integrate}. Which takes an expression and a variable.


\begin{api}
sympy.integrate(expression, variable)
\end{api}



For example to compute \(\int 4x^3 + 2x + 1 dx\):




\begin{pyin}
sym.integrate(expression, x)
\end{pyin}




\begin{equation*}
\begin{split}\displaystyle x^{4} + x^{2} + x\end{split}
\end{equation*}




\subsection{Calculate the definite integral of an expression.}

You can calculate the definite integral of an expression using
\texttt{sympy.integrate}. The first argument is an expression but instead of passing a
variable as the second argument you pass a tuple with the variable as well as the upper
and lower bounds of integration.


\begin{api}
sympy.integrate(expression, (variable, lower_bound, upper_bound))
\end{api}



For example to compute \(\int_0^4 4x^3 + 2x + 1 dx\):




\begin{pyin}
sym.integrate(expression, (x, 0, 4))
\end{pyin}




\begin{equation*}
\begin{split}\displaystyle 276\end{split}
\end{equation*}




\subsection{Use \(\infty\protect\)}
\label{\detokenize{tools-for-mathematics/03-calculus/how/main:use-infty}}

In sympy you can access \(\infty\) using \texttt{sym.oo}:


\begin{api}
sympy.oo
\end{api}



For example:




\begin{pyin}
sym.oo
\end{pyin}




\begin{equation*}
\begin{split}\displaystyle \infty\end{split}
\end{equation*}




\subsection{Calculate \index{limits}}

You can calculate \index{limits} using \texttt{sympy.limit}. The first argument is the
expression, then the variable and finally the expression the variable tends to.


\begin{api}
sympy.limit(expression, variable, value)
\end{api}



For example to compute \(\lim_{h \to 0} \frac{4 x ^ 3 + 2 x + 1 - 4(x - h)^3 - 2(x - h) - 1}{h}\):




\begin{pyin}
h = sym.Symbol("h")
expression = (4 * x ** 3 + 2 * x + 1 - 4 * (x - h) ** 3 - 2 * (x - h) - 1) / h
sym.limit(expression, h, 0)
\end{pyin}




\begin{equation*}
\begin{split}\displaystyle 12 x^{2} + 2\end{split}
\end{equation*}






\section{Exercises}
\label{\detokenize{tools-for-mathematics/03-calculus/exercises/main:exercises}}\label{\detokenize{tools-for-mathematics/03-calculus/exercises/main::doc}}
\begin{enumerate}

\item 

For each of the following functions calculate \(\frac{df}{dx}\), \(\frac{d^2f}{dx^2}\) and \(\int f(x) dx\).
\begin{enumerate}

\item 

\(f(x) = x\)

\item 

\(f(x) = x ^{\frac{1}{3}}\)

\item 

\(f(x) = 2 x (x - 3) (\sin(x) - 5)\)

\item 

\(f(x) = 3  x ^ 3 + 6 \sqrt{x} + 3\)

\end{enumerate}

\item 

Consider the function \(f(x)=2x+1\). By \index{differentiating} \textbf{from first principles} show that \(f'(x)=2\).

\item 

Consider the second derivative \(f''(x)=6x+4\) of some cubic function \(f(x)\).
\begin{enumerate}

\item 

Find \(f'(x)\)

\item 

You are given that \(f(0)=10\) and \(f(1)=13\), find \(f(x)\).

\item 

Find all the stationary points of \(f(x)\) and determine their nature.

\end{enumerate}

\item 

Consider the function \(f(x)=\frac{2}{3}x ^ 3 + b x ^ 2 + 2 x + 3\), where \(b\) is some undetermined coefficient.
\begin{enumerate}

\item 

Find \(f'(x)\) and \(f''(x)\)

\item 

You are given that \(f(x)\) has a stationary point at \(x=2\). Use this information to find \(b\).

\item 

Find the coordinates of the other stationary point.

\item 

Determine the nature of all stationary points.

\end{enumerate}

\item 

Consider the functions \(f(x)=-x^2+4x+4\) and \(g(x)=3x^2-2x-2\).
\begin{enumerate}

\item 

Create a variable \texttt{turning\_points} which has value the turning points of \(f(x)\).

\item 

Create variable \texttt{intersection\_points} which has value of the points where \(f(x)\) and \(g(x)\) intersect.

\item 

Using your ansyours to parts 2., calculate the area of the region betyouen \(f\) and \(g\). Assign this value to a variable \texttt{area\_betyouen}.

\end{enumerate}

\end{enumerate}

\section{Further information}
\label{\detokenize{tools-for-mathematics/03-calculus/why/main:further-information}}\label{\detokenize{tools-for-mathematics/03-calculus/why/main::doc}}

\subsection{How can you plot a function}
\label{\detokenize{tools-for-mathematics/03-calculus/why/main:how-can-you-plot-a-function}}

It is possible to plot a function using \texttt{sympy} using the \texttt{sympy.plot} function:


\begin{api}
sympy.plot(expression)
\end{api}



So for example, here is a plot of \(f(x)=x^2 + 3x + 1\):





\begin{pyin}
import sympy as sym

x = sym.Symbol("x")
sym.plot(x ** 2 + 3 * x + 1)
\end{pyin}


\begin{center}
    \includegraphics[width=.7\textwidth]{./assets/plotting_with_sympy/main.pdf}
\end{center}

It is possible to specify the \(x\) \index{limits} and combine it with other plots:

\begin{pyin}
sym.plot(x ** 2 + 3 * x + 1, xlim=(-5, 5))
\end{pyin}

\begin{center}
    \includegraphics[width=.7\textwidth]{./assets/plotting_with_sympy_with_xlim/main.pdf}
\end{center}




\textbf{This plotting solution is good if you want to take a look at a function
quickly but it is not recommended.} The main python library for plotting is called
\texttt{matplotlib} and is covered in a chapter of the online version of the
book.
